\documentclass[aspectratio=169,11pt]{beamer}
\makeatletter
\input{../../templates/dune-professional/beamerthemeDUNEProfessional.sty}
\makeatother
\title{DPS hardware for vertical drift}
\subtitle{Existing BDE design, additions and procurement}
\author{Manuel Arroyave}
\date{I\&I workshop draft --- 27 October 2026}
\newcommand{\mydate}{27 October 2026}
\hypersetup{colorlinks=true,urlcolor=DuneInfo,linkcolor=DuneInfo}
\newcommand{\planningnote}[1]{\par\medskip{\footnotesize\color{DuneTextMuted}#1\par}}
\begin{document}

\begin{frame}[plain]
  \titlepage
\end{frame}

\begin{frame}{Use a commercial PLC for VD detector protection}
  \large
  \textbf{Proposal:} use the bottom drift electronics (BDE) design
  and add the inputs and outputs needed by the other VD systems.\par
  \bigskip
  Put the central programmable logic controller (PLC) in the CUC,
  with distributed I/O near the equipment.\par
  \bigskip
  Use OPC UA to connect DPS status and agreed commands to slow controls.\par
  \planningnote{CUC = Central Utility Cavern; central-PLC location proposed. HD can later use an extension or a separate PLC. This plan covers VD.}
\end{frame}

\begin{frame}{The VD rack plan already reserves space for DPS}
  \renewcommand{\arraystretch}{1.4}
  \begin{tabularx}{\textwidth}{@{}>{\RaggedRight\arraybackslash}p{3.4cm}r>{\RaggedRight\arraybackslash}X@{}}
    \toprule
    \textbf{Rack type} & \textbf{Count} & \textbf{DPS equipment / connections}\\
    \midrule
    BDE auxiliary & 10 & 12U per rack; eight WIEC connections per rack\\
    BDE low voltage & 20 & One PL506 per rack; four WIECs per supply\\
    Bias high voltage & 10 & 12U per rack for the DPS interface to ISEG supplies\\
    \bottomrule
  \end{tabularx}\par
  \bigskip
  Use these locations and rack allocations to place the additional I/O.
  \planningnote{May 2025 VD rack workbook; June 2026 BDE layout. U is rack height. These are design allocations; the bias-HV sheet is a first estimate.}
\end{frame}

\begin{frame}{One BDE I/O group serves eight WIECs}
  \centering
  \begin{tikzpicture}[
    box/.style={fill=DunePaper,text=DuneOnAccent,align=center,
      text width=3.1cm,minimum height=1.05cm,inner sep=6pt},
    edge/.style={-{Stealth},draw=DuneSky,line width=1.2pt}]
    \node[box,fill=DuneSky] (io) at (-4,0) {BDE auxiliary rack\\remote I/O};
    \node[box] (lv) at (1.6,1.1) {2 LV racks\\2 PL506 supplies};
    \node[box,fill=DuneCoral] (wiec) at (1.6,-1.1) {8 WIECs\\8 PTCs + 48 WIBs};
    \draw[edge] (io.north east) -| (-1.7,1.1) -- node[above,font=\small] {inhibits} (lv.west);
    \draw[edge,<->] (io.south east) -| (-1.7,-1.1) -- node[below,font=\small] {EtherCAT} (wiec.west);
    \draw[edge] (lv.east) -- (4.35,1.1) -- node[right,font=\small] {48 V} (4.35,-1.1) -- (wiec.east);
  \end{tikzpicture}\par
  \medskip
  \DuneDiagramCaption{Ten groups: 20 PL506 supplies, 80 WIECs and 480 WIBs.}
  \medskip
  \raggedright The auxiliary rack also connects the fans, heaters and temperature sensors.
  \planningnote{WIEC = warm interface electronics crate; PTC = power and timing card; WIB = warm interface board. June 2026 layout; fan/heater wiring omitted here.}
\end{frame}

\begin{frame}{Panel drawings and the Ignition front end are available}
  \begin{columns}[T,onlytextwidth]
    \begin{column}{0.47\textwidth}
      \textbf{TPC protection panel designs}\par\medskip
      \colorbox{DunePaper}{\includegraphics[page=2,width=\dimexpr\linewidth-2\fboxsep\relax,height=3.35cm,keepaspectratio,
        trim=290bp 500bp 900bp 430bp,clip]{figures/tpc-panel-reference.pdf}}\par
      \medskip\DuneDiagramCaption{Central PLC, bias I/O and LV/fan/heater panels.}
    \end{column}
    \begin{column}{0.49\textwidth}
      \textbf{BDE Ignition screens}\par\medskip
      \includegraphics[page=7,width=\linewidth,height=3.35cm,keepaspectratio,
        trim=10bp 210bp 460bp 74bp,clip]{figures/sc-progress-june2026.pdf}\par
      \planningnote{Existing display of WIEC groups; June 2026 SC progress presentation.}
    \end{column}
  \end{columns}\par
  \medskip
  Parts lists and power calculations are available too. Use them to specify the VD additions.
  \planningnote{TPC reference: EDMS 2401090 v4, sheet 2. Use the VD rack allocations when adapting these panels.}
\end{frame}

\begin{frame}{The PX3 rack-protection setup already defines shutdown}
  \renewcommand{\arraystretch}{1.45}
  \begin{tabularx}{\textwidth}{@{}>{\RaggedRight\arraybackslash}p{3.7cm}>{\RaggedRight\arraybackslash}X@{}}
    \toprule
    \textbf{June 2026 setup guide} & \textbf{Documented behaviour}\\
    \midrule
    Ignition connected & Ignition reads sensors and switches outlets over SNMP\\
    Ignition unreachable & PX3 local rules handle smoke, water and temperature; the network-switch outlet stays on\\
    After a fallback trip & Power is restored manually or through Ignition\\
    \bottomrule
  \end{tabularx}\par
  \bigskip
  Agree where the VD rack-protection logic will run: PDU, DPS PLC and Ignition each need a defined role.
  \planningnote{DUNE-doc-32655 v8, PX3 Rack Protection Setup, pp.~3--6 and 14--16. Confirm the selected rack/PDU configuration against this reference setup.}
\end{frame}

\begin{frame}{Use the inventory to specify the additional I/O}
  \renewcommand{\arraystretch}{1.35}
  \begin{tabularx}{\textwidth}{@{}>{\RaggedRight\arraybackslash}p{3cm}>{\RaggedRight\arraybackslash}X@{}}
    \toprule
    \textbf{Existing material} & \textbf{Work for VD}\\
    \midrule
    Rack and panel inventory & Allocate terminals, power and cable routes for added signals\\
    Subsystem hardware lists & Remove overlaps; complete CalCI and other missing connections\\
    Interface documents and matrix & Resolve conflicting actions; record signal type, timing and feedback\\
    \bottomrule
  \end{tabularx}\par
  \bigskip
  \textbf{Example:} 80 PoF boxes $\times$ (8 temperatures + 8 voltages)
  = 1,280 monitoring values.\par
  \planningnote{Monitoring values are not automatically PLC terminals. Choose the readout and protection interfaces before counting modules. CalCI = calibration and cryogenic instrumentation; PoF = power over fibre.}
\end{frame}

\begin{frame}{Confirm the time allowed for each protective action}
  \large
  So far, I have not found a critical protection function
  that needs a response below 1 ms.\par
  \bigskip
  Each consortium must specify the time from a hazardous condition
  to the completed protective action.\par
  \bigskip
  Include detection, communication, PLC processing and equipment shutdown.\par
  \planningnote{Use the simulation framework to prepare tests. Measure complete paths on hardware; add faster local protection wherever a confirmed requirement needs it.}
\end{frame}

\begin{frame}{Resolve the interfaces that affect hardware}
  \renewcommand{\arraystretch}{1.35}
  \begin{tabularx}{\textwidth}{@{}>{\RaggedRight\arraybackslash}p{3.3cm}>{\RaggedRight\arraybackslash}X@{}}
    \toprule
    \textbf{Interface} & \textbf{Decision needed}\\
    \midrule
    SC / DPS & With Geoff, agree command priority, reset and re-enable; use SC for relevant long-term records\\
    Cryogenic measurements & Confirm each source and whether DPS needs a dedicated signal or a software exchange\\
    LBNF/DUNE and SDSTA BMS & Select the limited hardware signals and their connection point; define software monitoring\\
    CUC / detector mezzanine & Specify fieldbus routes, networks and control-power supplies\\
    \bottomrule
  \end{tabularx}
  \planningnote{BMS = building management system. Keep the two BMSs distinct. SDSTA retains responsibility for personnel-safety alarms.}
\end{frame}

\begin{frame}{Use the SURF inspection to map the facility connections}
  \begin{columns}[T,onlytextwidth]
    \begin{column}{0.48\textwidth}
      \DuneFixedColorBlock{paper}{2.7cm}{1}{Already prepared}
        {Marked maps for levels 4850 and 4910, with drawing references and a panel survey register.}
    \end{column}
    \begin{column}{0.48\textwidth}
      \DuneFixedColorBlock{sky}{2.7cm}{2}{Record on site}
        {Controller and I/O models, panel labels, available terminals and connection routes underground and at Ross Dry.}
    \end{column}
  \end{columns}\par
  \bigskip
  Coordinate access with Rachel; ask Dale Curran about the SDSTA interfaces.
  \planningnote{Inspection planned for this week, subject to access and safety arrangements. Maps are prepared but have not yet been shared.}
\end{frame}

\begin{frame}{Use the existing test stands for the VD additions}
  \renewcommand{\arraystretch}{1.4}
  \begin{tabularx}{\textwidth}{@{}>{\RaggedRight\arraybackslash}X>{\RaggedRight\arraybackslash}X@{}}
    \toprule
    \textbf{Test} & \textbf{Measure or demonstrate}\\
    \midrule
    PLC, remote I/O, PTC and supply inhibit & Complete trip time and equipment feedback\\
    PDU sensors and shutdown logic & Response with SC connected, lost or stalled\\
    SC commands, reset and re-enable & Protection takes priority; recovery is deliberate\\
    Fieldbus, network and power failures & Output states, lost branches and restart behaviour\\
    \bottomrule
  \end{tabularx}
  \planningnote{Fermilab BDE/VST and G271 material already exists. Agree test time, configuration and extra equipment with the owners; identify any SURF test needed.}
\end{frame}

\begin{frame}{Prepare the VD design and purchases for June 2027}
  \renewcommand{\arraystretch}{1.4}
  \begin{tabularx}{\textwidth}{@{}>{\RaggedRight\arraybackslash}p{3.2cm}>{\RaggedRight\arraybackslash}X@{}}
    \toprule
    \textbf{Step} & \textbf{Result}\\
    \midrule
    Reconcile existing records & One VD rack, panel and I/O schedule; list added signals and unresolved requirements\\
    Design and prototype & Updated drawings, power budget and test results for the added interfaces\\
    June 2027 review & Parts list, costs, lead times, spares, installation sequence and procurement readiness\\
    After review & Order approved hardware; install against confirmed site dates\\
    \bottomrule
  \end{tabularx}
  \planningnote{Check possible summer 2027 facility-interface needs now. Full VD response-matrix review targets early 2028; hardware requirements are needed before purchasing.}
\end{frame}

\begin{frame}{Agree who does the work}
  \renewcommand{\arraystretch}{1.35}
  \begin{tabularx}{\textwidth}{@{}>{\RaggedRight\arraybackslash}p{4.1cm}>{\RaggedRight\arraybackslash}X@{}}
    \toprule
    \textbf{Proposed team} & \textbf{Next action}\\
    \midrule
    Manuel + control engineer & Specify VD additions, cost and procurement work\\
    Terri, Trevor + Manuel & Reuse the BDE hardware design and test arrangements\\
    Manuel + consortia & Resolve missing signals, timing and conflicting requirements\\
    Manuel + Geoff & Agree SC/DPS commands, rack protection and recovery\\
    Vivian + subsystem / DPS teams & Update ICDs and record the agreed interfaces\\
    \bottomrule
  \end{tabularx}
  \planningnote{Confirm support and dates at I\&I. Agree who protects equipment in each installation and operating phase before DPS is available.}
\end{frame}

\appendix
\begin{frame}{Reference: the May 2025 VD rack inventory}
  \renewcommand{\arraystretch}{1.15}
  \begin{tabularx}{\textwidth}{@{}>{\RaggedRight\arraybackslash}Xr>{\RaggedRight\arraybackslash}Xr@{}}
    \toprule
    \textbf{Rack type} & \textbf{Count} & \textbf{Rack type} & \textbf{Count}\\
    \midrule
    BDE low voltage & 20 & PDS low voltage & 4\\
    BDE auxiliary & 10 & Drift HV & 2\\
    TDE & 20 & Reserved drift HV & 2\\
    Bias HV & 10 & Spare / CalCI & 8\\
    \bottomrule
  \end{tabularx}\par
  \bigskip
  \textbf{76 numbered mezzanine positions:} SEDM-01 to SEDM-76.\par
  \medskip
  The worksheet also has a separate concept for 40 top-plate mini-racks.
  \planningnote{Counts are rack positions, not numbers of PLC panels or remote-I/O stations. Reconcile this May allocation with later changes before ordering.}
\end{frame}
\end{document}
